%Data institusi

%\input{Dozentdaten}

%\renewcommand{\fachbereich}{Bidang studi}

%\renewcommand{\dozent}{Prof. Dr. . }

%Data ujian

\renewcommand{\klausurgebiet}{ }

\renewcommand{\klausurtyp}{ }

\renewcommand{\klausurdatum}{. 20}

\klausurvorspann {\fachbereich} {\klausurdatum} {\dozent} {\klausurgebiet} {\klausurtyp}

%Data untuk tabel poin berikut

\renewcommand{\aeins}{  3  }

\renewcommand{\azwei}{  3   }

\renewcommand{\adrei}{  6  }

\renewcommand{\avier}{  4  }

\renewcommand{\afuenf}{  5  }

\renewcommand{\asechs}{  0  }

\renewcommand{\asieben}{  5  }

\renewcommand{\aacht}{  6  }

\renewcommand{\aneun}{  6  }

\renewcommand{\azehn}{  0  }

\renewcommand{\aelf}{  6  }

\renewcommand{\azwoelf}{  12  }

\renewcommand{\adreizehn}{  0   }

\renewcommand{\avierzehn}{  9  }

\renewcommand{\afuenfzehn}{  65  }

\renewcommand{\asechzehn}{    }

\renewcommand{\asiebzehn}{    }

\renewcommand{\aachtzehn}{    }

\renewcommand{\aneunzehn}{    }

\renewcommand{\azwanzig}{    }

\renewcommand{\aeinundzwanzig}{    }

\renewcommand{\azweiundzwanzig}{    }

\renewcommand{\adreiundzwanzig}{    }

\renewcommand{\avierundzwanzig}{    }

\renewcommand{\afuenfundzwanzig}{    }

\renewcommand{\asechsundzwanzig}{    }

\punktetabellevierzehn

\klausurnote

\newpage

\setcounter{section}{0}



\inputaufgabeklausurloesung
{3}

{

Definisikan istilah-istilah berikut
\zusatzklammer {yang dicetak miring} {} {}.
\aufzaehlungsechs{Suatu
\stichwort {kelengkungan utama} {}
pada suatu permukaan diferensiabel

\mavergleichskette
{\vergleichskette
{ Y
}
{ \subseteq }{ \R^3
}
{  }{
}
{    }{
}
{   }{
}
}
{}{}{}
pada titik

\mavergleichskette
{\vergleichskette
{ P
}
{ \in }{ Y
}
{  }{
}
{    }{
}
{   }{
}
}
{}{}{.}

}{Suatu
\stichwort {difeomorfisme} {}
antara manifold diferensiabel
\mathkor {} {L} {dan} {M} {.}

}{Suatu
\stichwort {vektor singgung} {}
di titik

\mavergleichskette
{\vergleichskette
{  P
}
{ \in }{  M
}
{  }{
}
{    }{
}
{   }{
}
}
{}{}{}
dari manifold diferensiabel $M$.

}{Suatu
\stichwort {operator diferensial orde pertama} {}
pada manifold diferensiabel $M$.

}{Suatu
\stichwort {setengah ruang Euklides} {.}

}{Suatu
\definitionsverweis {penghabisan kompak}{}{}
bagi suatu
\definitionsverweis {ruang topologis}{}{}
$X$.
}

}
{

\aufzaehlungsechs{Setiap
\definitionsverweis {nilai eigen}{}{}
dari
\definitionsverweis {pemetaan Weingarten}{}{}
\maabbeledisp {L_P} {T_P Y} { T_PY
} {v} {- \left(  D_{v} N  \right) \left(  P  \right)
} {,}
disebut
\definitionswort {kelengkungan utama}{}
dari $Y$ di $P$.
}{Suatu
\definitionsverweis {homeomorfisme}{}{}
\maabbdisp {\varphi} {L} {M
} {}
disebut
\stichwortpraemath {C^1} {difeomorfisme}{,}
jika
\mathkor {} {\varphi} {maupun} {\varphi^{-1}} {}
merupakan
$C^1$-\definitionsverweis {pemetaan}{}{}.
}{Yang dimaksud dengan vektor tangen di $P$ adalah suatu
\definitionsverweis {kelas ekuivalensi}{}{}
dari
\definitionsverweis {kurva-kurva diferensiabel}{}{}
yang
\definitionsverweis {ekuivalen secara tangensial}{}{}
melalui $P$.
}{Suatu pemetaan
\maabbdisp {D} {C^1(M,\R) } {C^0(M,\R)
} {}
disebut
operator diferensial orde pertama
jika memenuhi sifat-sifat berikut.
\aufzaehlungzwei {$D$ bersifat
$\R$-\definitionsverweis {linear}{}{.}
} {Berlaku

\mavergleichskette
{\vergleichskette
{ D(fg)
}
{ = }{ fD(g)+gD(f)
}
{  }{
}
{    }{
}
{   }{
}
}
{}{}{.}
}
}{Yang dimaksud dengan setengah ruang Euklides berdimensi $n$ adalah himpunan

\mavergleichskettedisp
{\vergleichskette
{ H
}
{ =} { \left\{ x \in \R^n   \mid  x_1 \geq 0        \right\}
}
{ } {
}
{ } {
}
{ } {
}
}
{}{}{}
dengan
\definitionsverweis {topologi terinduksi}{}{.}
}{Suatu penghabisan kompak

\mathbed {A_n} {}
 {n \in \N} {}
 {} {} {} {,}
dari $X$ adalah suatu
\definitionsverweis {barisan}{}{}
dari
\definitionsverweis {subhimpunan-subhimpunan kompak}{}{}

\mavergleichskette
{\vergleichskette
{ A_n
}
{ \subseteq }{ X
}
{  }{
}
{    }{
}
{   }{
}
}
{}{}{}
dengan

\mathdisp {A_n \subseteq A_{n+1}^{o} \text{ dan } \bigcup_{n=0}^\infty A_n = X} { . }

}

}



\inputaufgabeklausurloesung
{3}

{

Nyatakan teorema-teorema berikut.
\aufzaehlungdrei{Teorema tentang
\definitionsverweis {ruang singgung}{}{}
dari suatu
\definitionsverweis {submanifold tertutup}{}{}

\mavergleichskette
{\vergleichskette
{  M
}
{ \subseteq }{  N
}
{  }{
}
{    }{
}
{   }{
}
}
{}{}{.}}{Teorema tentang bentuk eksplisit dari bentuk diferensial tarik balik $\varphi^* \omega$ oleh pemetaan diferensiabel
\maabbdisp {\varphi} { U } { V
} {}
\zusatzklammer {dengan
\definitionsverweis {himpunan bagian terbuka}{}{}
\mathkor {} {U \subseteq \R^n} {dan} {V \subseteq \R^m} {,}
yang koordinatnya masing-masing dinyatakan dengan
\mathl{x_1  , \ldots , x_n}{} dan
\mathl{y_1  , \ldots , y_m}{}} {} {}
dari suatu
$k$-\definitionsverweis {bentuk diferensial}{}{}
pada $V$ dengan representasi

\mavergleichskettedisp
{\vergleichskette
{  \omega
}
{ =} {\sum_{I,\,  { \# \left(   I  \right) } = k}  f_I  dy_I
}
{ } {
}
{ } {
}
{ } {
}
}
{}{}{,}
dengan
\maabb {f_I} { V } {\R
} {}
adalah
\definitionsverweis {fungsi}{}{.}}{Teorema tentang karakterisasi kurva geodesik terhadap koneksi Levi-Civita pada manifold Riemann $M$.}

}
{

\aufzaehlungdrei{\faktsituation {Misalkan $N$ suatu
\definitionsverweis {manifold diferensiabel}{}{}
ber-
\definitionsverweis {dimensi}{}{}
$n$ dan misalkan

\mavergleichskette
{\vergleichskette
{ M
}
{ \subseteq }{ N
}
{  }{
}
{    }{
}
{   }{
}
}
{}{}{}
suatu
\definitionsverweis {submanifold tertutup}{}{}
berdimensi $m$.}
\faktfolgerung {Maka untuk setiap titik

\mavergleichskette
{\vergleichskette
{ P
}
{ \in }{ M
}
{  }{
}
{    }{
}
{   }{
}
}
{}{}{}
\definitionsverweis {pemetaan tangen}{}{}
\maabbdisp {} { T_PM } { T_PN
} {}
bersifat
\definitionsverweis {injektif}{}{.}}
\faktzusatz {Artinya, ruang tangen
\mathl{T_PM}{} merupakan
\definitionsverweis {subruang vektor}{}{}
berdimensi $m$ dari
\mathl{T_PN}{.}}
\faktzusatz {}}{\faktsituation {Misalkan
\mathkor {} {U \subseteq \R^n} {dan} {V \subseteq \R^m} {}
merupakan
\definitionsverweis {subhimpunan-subhimpunan terbuka}{}{,}
yang koordinatnya masing-masing dinyatakan dengan
\mathl{x_1 , \ldots , x_n}{} dan
\mathl{y_1 , \ldots , y_m}{}. Misalkan
\maabbdisp {\varphi} { U } { V
} {}
suatu
\definitionsverweis {pemetaan diferensiabel}{}{}
dan misalkan $\omega$ suatu
$k$-\definitionsverweis {bentuk diferensial}{}{}
pada $V$ dengan representasi

\mavergleichskettedisp
{\vergleichskette
{ \omega
}
{ =} { \sum_{I,\, { \# \left(   I  \right) } = k}  f_I  dy_I
}
{ } {
}
{ } {
}
{ } {
}
}
{}{}{,}
dengan
\maabb {f_I} { V } {\R
} {}
adalah
\definitionsverweis {
fungsi-fungsi}{}{.}}
\faktfolgerung {Maka
\definitionsverweis {bentuk tarikan balik}{}{}
memiliki representasi

\mavergleichskettealignhandlinks
{\vergleichskettealignhandlinks
{ \varphi^* \omega
}
{ =} { \sum_{I,\, { \# \left(   I  \right) } = k}  \left(  f_I \circ \varphi  \right)  \left( \sum_{J,\, { \# \left(   J  \right) } = k} \det  \left(  \left(  { \frac{   \partial  \varphi_i   }{  \partial   x_j   } }  \right)_{ i \in I, j \in J}  \right)  dx_J \right)
}
{ =} { \sum_{J,\, { \# \left(   J  \right) } = k}   \left( \sum_{I,\, { \# \left(   I  \right) } = k} \left(  f_I \circ \varphi  \right)  \det  \left(  \left(  { \frac{   \partial  \varphi_i   }{  \partial   x_j   } }  \right)_{ i \in I, j \in J}  \right)  \right)  dx_J
}
{ } {
}
{ } {
}
}
{}
{}{.}}
\faktzusatz {}
\faktzusatz {}}{\faktsituation {Suatu kurva yang dua kali
\definitionsverweis {diferensiabel}{}{}
\maabb {\gamma} { I } { M
} {}
pada suatu
\definitionsverweis {manifold Riemann}{}{}
$M$ adalah}
\faktfolgerung {suatu
\definitionsverweis {kurva geodesik}{}{}
\zusatzklammer {terhadap
\definitionsverweis {koneksi Levi-Civita}{}{}} {} {,}
tepat ketika dalam setiap peta kurva tersebut memenuhi
\definitionsverweis {sistem persamaan diferensial biasa orde dua}{}{}

\mavergleichskettedisp
{\vergleichskette
{ \gamma_k^{\prime \prime} + \sum_{ij} \Gamma^k_{ij} \gamma_i' \gamma_j'
}
{ =} { 0
}
{ } {
}
{ } {
}
{ } {
}
}
{}{}{}
\zusatzklammer {untuk semua $k$} {} {}
yang diberikan.}
\faktzusatz {}
\faktzusatz {}}

}



\inputaufgabeklausurloesung
{6 (2+4)}

{

Misalkan

\mavergleichskette
{\vergleichskette
{ \alpha,\beta
}
{ \in }{ \R_+
}
{  }{
}
{    }{
}
{   }{
}
}
{}{}{}
dan misalkan

\mavergleichskettedisp
{\vergleichskette
{ Y
}
{ =} { \left\{ (x,y)\in\R^2  \mid   \alpha x^2+\beta y^2 = 1        \right\}
}
{ } {
}
{ } {
}
{ } {
}
}
{}{}{.}
\aufzaehlungzwei {Tentukan
\definitionsverweis {pemetaan Gauss}{}{}
untuk $Y$.
} { Berikan secara eksplisit suatu
\definitionsverweis {pemetaan invers}{}{}
dari pemetaan Gauss untuk $Y$.
}

}
{

\aufzaehlungzwei {Sebagai
\definitionsverweis {medan normal satuan}{}{},
kita dapat mulai dari gradien fungsi pendefinisi, yaitu

\mavergleichskettedisp
{\vergleichskette
{ G(x,y)
}
{ =} { \begin{pmatrix}  2 \alpha x  \\ 2 \beta y    \end{pmatrix}
}
{ } {
}
{ } {
}
{ } {
}
}
{}{}{}
dan setelah normalisasi diperoleh

\mavergleichskettedisp
{\vergleichskette
{ N(x,y)
}
{ =} {  { \frac{   1 }{  \sqrt{ \alpha^2 x^2 + \beta^2 y^2 }  } } \begin{pmatrix}   \alpha x  \\  \beta y    \end{pmatrix}
}
{ } {
}
{ } {
}
{ } {
}
}
{}{}{.}
Jadi, pemetaan Gauss adalah
\maabbeledisp {\varphi} {Y = \left\{  \begin{pmatrix}  x  \\ y   \end{pmatrix}   \in \R^2  \mid   \alpha x^2+\beta y^2 = 1        \right\}  } { S^1
} { \begin{pmatrix}  x  \\ y   \end{pmatrix} } { { \frac{   1 }{  \sqrt{ \alpha^2 x^2 + \beta^2 y^2 }  } } \begin{pmatrix}   \alpha x  \\  \beta y    \end{pmatrix}
} {.}
} {Kita nyatakan bahwa
\maabbeledisp {\psi} {S^1} {Y
} { \begin{pmatrix}  u  \\v   \end{pmatrix} } { { \frac{   1 }{  \sqrt{ { \frac{  u^2 }{ \alpha } }   + { \frac{  v^2 }{ \beta } } }  } }       \begin{pmatrix}   { \frac{   u  }{ \alpha } }   \\ { \frac{   v  }{ \beta } }    \end{pmatrix}
} {,}
adalah fungsi invers. Karena

\mavergleichskettedisp
{\vergleichskette
{ { \frac{   1 }{   { \frac{  u^2 }{ \alpha } }   + { \frac{  v^2 }{ \beta } }  } } \left(  \alpha { \frac{  u^2 }{ \alpha^2 } } + \beta{ \frac{  v^2 }{ \beta^2 } }  \right)
}
{ =} { 1
}
{ } {
}
{ } {
}
{ } {
}
}
{}{}{}
citra $\psi$ terletak dalam $Y$. Karena

\mavergleichskettealign
{\vergleichskettealign
{  \varphi \left(  \psi \begin{pmatrix}  u  \\v   \end{pmatrix}  \right)
}
{ =} { \varphi \left(  { \frac{   1 }{  \sqrt{ { \frac{  u^2 }{ \alpha } }   + { \frac{  v^2 }{ \beta } } }  } }       \begin{pmatrix}   { \frac{   u  }{ \alpha } }   \\ { \frac{   v  }{ \beta } }    \end{pmatrix}    \right)
}
{ =} { \varphi \begin{pmatrix}   { \frac{   1 }{  \alpha    \sqrt{ { \frac{  u^2 }{ \alpha } }   + { \frac{  v^2 }{ \beta } } }  } }   u    \\ { \frac{   1 }{ \beta  \sqrt{ { \frac{  u^2 }{ \alpha } }  + { \frac{  v^2 }{ \beta } } }  } }  v       \end{pmatrix}
}
{ =} { { \frac{   1 }{  \sqrt{  \alpha^2 { \frac{  u^2 }{  \alpha^2   \left(   { \frac{  u^2 }{ \alpha } }  +  { \frac{  v^2 }{ \beta } }   \right)  } }  + \beta^2 { \frac{  v^2 }{  \beta^2 \left(   { \frac{  u^2 }{ \alpha } }  +  { \frac{  v^2 }{ \beta } }   \right)  } }       }  } }  \begin{pmatrix}  { \frac{   1 }{  \sqrt{ { \frac{  u^2 }{ \alpha } }   + { \frac{  v^2 }{ \beta } } }  } }   u    \\   { \frac{   1 }{   \sqrt{ { \frac{  u^2 }{ \alpha } } + { \frac{  v^2 }{ \beta } } }  } } v    \end{pmatrix}
}
{ =} { { \frac{   1 }{  \sqrt{ { \frac{  u^2 }{  \left(   { \frac{  u^2 }{ \alpha } }  +  { \frac{  v^2 }{ \beta } }   \right)  } }   + { \frac{  v^2 }{  \left(   { \frac{  u^2 }{ \alpha } }  +  { \frac{  v^2 }{ \beta } }   \right)  } }       }  } }  \begin{pmatrix}   { \frac{   1 }{  \sqrt{ { \frac{  u^2 }{ \alpha } }   + { \frac{  v^2 }{ \beta } } }  } }   u    \\   { \frac{   1 }{   \sqrt{ { \frac{  u^2 }{ \alpha } }   + { \frac{  v^2 }{ \beta } } }  } } v    \end{pmatrix}
}
}
{
\vergleichskettefortsetzungalign
{ =} { { \frac{    \sqrt{ { \frac{  u^2 }{ \alpha } }  +  { \frac{  v^2 }{ \beta } } } }{  \sqrt{    u^2+ v^2       }  } }  \begin{pmatrix}   { \frac{   1 }{  \sqrt{ { \frac{  u^2 }{ \alpha } }   + { \frac{  v^2 }{ \beta } } }  } }   u    \\   { \frac{   1 }{   \sqrt{ { \frac{  u^2 }{ \alpha } }   + { \frac{  v^2 }{ \beta } } }  } } v    \end{pmatrix}
}
{ =} { \begin{pmatrix}  u  \\v   \end{pmatrix}
}
{ } {
}
{ } {}
}
{}{}
dan

\mavergleichskettealign
{\vergleichskettealign
{  \psi \left(  \varphi \begin{pmatrix}  x  \\ y   \end{pmatrix}  \right)
}
{ =} { \psi \begin{pmatrix}  { \frac{    \alpha x }{  \sqrt{\alpha^2 x^2+ \beta^2 y^2 }  } }   \\ { \frac{   \beta y }{  \sqrt{\alpha^2 x^2+ \beta^2 y^2  } }}    \end{pmatrix}
}
{ =} { { \frac{   1 }{  \sqrt{ { \frac{   \alpha x^2 }{  \alpha^2x^2 + \beta^2y^2  } }  + { \frac{   \beta y^2 }{  \alpha^2x^2 + \beta^2y^2  } }  }  } }  \begin{pmatrix}  { \frac{    x }{  \sqrt{\alpha^2 x^2+ \beta^2 y^2 }  } }   \\ { \frac{   y }{  \sqrt{\alpha^2 x^2+ \beta^2 y^2  } }}    \end{pmatrix}
}
{ =} { { \frac{   \sqrt{\alpha^2 x^2+ \beta^2 y^2 } }{  \sqrt{    \alpha x^2 +  \beta y^2  }  } }  \begin{pmatrix}  { \frac{    x }{  \sqrt{\alpha^2 x^2+ \beta^2 y^2 }  } }   \\ { \frac{   y }{  \sqrt{\alpha^2 x^2+ \beta^2 y^2  } }}    \end{pmatrix}
}
{ =} { \begin{pmatrix}  x  \\ y   \end{pmatrix}
}
}
{}
{}{}
kedua pemetaan tersebut saling invers.
}

}



\inputaufgabeklausurloesung
{4}

{

Jelaskan pasangan istilah \anfuehrung{ekstrinsik dan intrinsik}{} dalam konteks manifold.

}
{  }



\inputaufgabeklausurloesung
{5}

{

Buktikan rumus luas permukaan suatu benda putar yang dihasilkan oleh
\definitionsverweis {kurva diferensiabel}{}{}
\maabbeledisp {\gamma} {]a,b[} {\R^2
} { t } {(x(t),y(t))
} {,}
dengan

\mavergleichskette
{\vergleichskette
{ y(t)
}
{ > }{ 0
}
{  }{
}
{    }{
}
{   }{
}
}
{}{}{.}

}
{

Misalkan $S$ permukaan putar, yang merupakan submanifold tertutup berdimensi dua dalam suatu himpunan terbuka di $\R^3$. Kita terapkan
Korolar 17.2 (Geometri diferensial (Osnabrück 2023))
pada parametrisasi
\maabbeledisp {} { ]a,b[ \times ]0, 2 \pi[} {S
} { (t, \alpha) } { (x(t),  y(t) \cos \alpha  , y(t) \sin \alpha  )
} {,}.
Turunan-turunan parsialnya adalah

\mathdisp {\begin{pmatrix}  x'(t) \\ y'(t) \cos \alpha \\  y'(t) \sin \alpha   \end{pmatrix} \text{   dan } \begin{pmatrix}  0  \\ - y(t) \sin \alpha \\  y(t) \cos \alpha   \end{pmatrix}} {  }

dan karena itu

\mathdisp {E(t, \alpha) = (x'(t))^2 + (y'(t))^2,\, F(t, \alpha)=0,\, G(t, \alpha) = (y(t))^2} { . }

Jadi, luas permukaannya sama dengan

\mavergleichskettedisphandlinks
{\vergleichskettedisphandlinks
{ \int_{  a  }^{  b  } \int_{  0  }^{  2 \pi } \sqrt{(x'(t))^2+(y'(t))^2} \cdot   y(t) \, d \alpha \, d t
}
{ =} { 2 \pi \int_{  a  }^{  b  } \sqrt{(x'(t))^2+(y'(t))^2} \cdot   y(t) \, d t
}
{ } {
}
{ } {
}
{ } {
}
}
{}{}{.}

}



\inputaufgabeklausurloesung
{0}

{

}
{  }



\inputaufgabeklausurloesung
{5}

{

Misalkan $X$ adalah suatu
\definitionsverweis {torus}{}{.}
Berikan suatu
\definitionsverweis {pemetaan diferensiabel}{}{}
surjektif
\maabbdisp {\varphi} {\R^2} {X
} {}
sedemikian sehingga
\definitionsverweis {pemetaan singgung}{}{}
\maabbdisp {T_P(\varphi)} { T_P\R^2 } { T_{\varphi(P)}X
} {}
juga surjektif di setiap titik

\mavergleichskette
{\vergleichskette
{ P
}
{ \in }{ \R^2
}
{  }{
}
{    }{
}
{   }{
}
}
{}{}{.}

}
{

Kita tulis torus sebagai $X=S^1 \times S^1$ dengan lingkaran satuan
\mathl{S^1= \left\{ (x,y) \in \R^2  \mid   \betrag { (x,y) } = 1        \right\}}{.}
Pemetaan
\maabbeledisp {\varphi} {\R^2} {S^1 \times S^1
} {(u,v)} {( \cos  u, \sin  u , \cos  v , \sin  v )
} {,}
bersifat surjektif karena kedua komponennya merupakan parametrisasi standar lingkaran satuan. Surjektivitas pemetaan tangen juga dapat diperiksa per komponen karena ruang tangen manifold produk adalah hasil kali ruang-ruang tangennya. Turunan dari
\maabbeledisp {} {\R} { \R^2
} { t } { ( \cos  t, \sin  t)
} {,}
adalah
\mathl{(- \sin  t, \cos  t) \neq (0,0)}{,} sehingga vektor citra ini selalu merentang ruang tangen satu-dimensional dari lingkaran satuan di titik citra.

}



\inputaufgabeklausurloesung
{6 (2+2+2)}

{

Tinjau grafik $M$ dari pemetaan
\maabbeledisp {\varphi} {\R^2} {\R
} {(u,v)} { u^2+uv-v^3
} {,}
sebagai submanifold tertutup berdimensi dua di $\R^3$, yaitu

\mavergleichskettedisp
{\vergleichskette
{ M
}
{ =} { \left\{ (u,v,u^2+uv-v^3)  \mid  (u,v) \in \R^2         \right\}
}
{ } {
}
{ } {
}
{ } {
}
}
{}{}{}
dengan metrik Riemann yang diinduksi dari $\R^3$. Misalkan
\maabbeledisp {\psi} {\R^2} { M
} {(u,v)} { (u,v,u^2+uv-v^3)
} {,}
adalah difeomorfisme yang bersesuaian.

a) Tentukan diferensial total dari $\psi$ serta vektor-vektor citra
\mathl{T_P(\psi) (e_1)}{} dan
\mathl{T_P(\psi) (e_2)}{} di
\mathl{T_{\psi(P)}M}{.}

b) Untuk setiap titik berbentuk

\mavergleichskette
{\vergleichskette
{ P
}
{ = }{ (u,0)
}
{  }{
}
{    }{
}
{   }{
}
}
{}{}{}
tentukan luas jajaran genjang yang direntang oleh
\mathl{T_P(\psi) (e_1)}{} dan
\mathl{T_P(\psi) (e_2)}{} di
\mathl{T_{\psi(P)}M}{}.

c) Untuk setiap titik berbentuk

\mavergleichskette
{\vergleichskette
{ P
}
{ = }{ (0,v)
}
{  }{
}
{    }{
}
{   }{
}
}
{}{}{}
tentukan luas jajaran genjang yang direntang oleh
\mathl{T_P(\psi) (e_1)}{} dan
\mathl{T_P(\psi) (e_2)}{} di
\mathl{T_{\psi(P)}M}{}.

}
{

a) Diferensial total dari $\psi$ di titik
\mathl{P=(u,v)}{} adalah

\mathdisp {\left(D \psi\right)_{P} = \begin{pmatrix}  1 &  0 \\  0 &  1 \\   2u+v & u-3v^2 \end{pmatrix}} {  }

dan berlaku

\mathdisp {T_P(\psi) (e_1) = \left(  D \psi  \right)_{ P } \left(  e_1  \right) = \begin{pmatrix}  1 \\ 0\\  2u+v   \end{pmatrix} \text{   dan } T_P(\psi) (e_2)   = \left(  D \psi  \right)_{ P } \left(  e_2  \right) = \begin{pmatrix}  0 \\ 1\\ u-3v^2   \end{pmatrix}} { . }

b) dan c) Untuk menentukan luas, pertama-tama kita hitung hasil kali skalar dari kedua vektor
\mathkor {} {\begin{pmatrix}  1 \\ 0\\  2u+v   \end{pmatrix}} {dan} {\begin{pmatrix}  0 \\ 1\\ u-3v^2   \end{pmatrix}} {.}
Kita peroleh

\mathdisp {\left\langle \begin{pmatrix}  1 \\ 0\\  2u+v   \end{pmatrix}  , \begin{pmatrix}  1 \\ 0\\  2u+v   \end{pmatrix}   \right\rangle = 1+(2u+v)^2 = 1 +4u^2+4uv +v^2} { , }

\mathdisp {\left\langle \begin{pmatrix}  1 \\ 0\\  2u+v   \end{pmatrix}  , \begin{pmatrix}  0 \\ 1\\ u-3v^2   \end{pmatrix}    \right\rangle = (2u+v)(u-3v^2) = 2u^2 -6uv^2 +uv -3v^3} {  }

dan

\mathdisp {\left\langle \begin{pmatrix}  0 \\ 1\\ u-3v^2   \end{pmatrix}   , \begin{pmatrix}  0 \\ 1\\ u-3v^2   \end{pmatrix}    \right\rangle = 1+(u-3v^2)^2 = 1 +u^2 -6uv^2  + 9v^4} { . }

b) Untuk
\mathl{P=(u,0)}{}, luas jajaran genjang yang direntang oleh
\mathl{T_P(\psi) (e_1)}{} dan
\mathl{T_P(\psi) (e_2)}{} di
\mathl{T_{\psi(P)}M}{} adalah

\mavergleichskettedisp
{\vergleichskette
{ \sqrt { \det  \begin{pmatrix}  1+4u^2  &   2u^2   \\  2 u^2 &  1+ u^2  \end{pmatrix} }
}
{ =} { \sqrt{  (1+4u^2)(1+u^2) - 4u^4 }
}
{ =} { \sqrt{  1+5u^2 }
}
{ } {
}
{ } {
}
}
{}{}{.}

c) Untuk
\mathl{P=(0,v)}{}, luas jajaran genjang yang direntang oleh
\mathl{T_P(\psi) (e_1)}{} dan
\mathl{T_P(\psi) (e_2)}{} di
\mathl{T_{\psi(P)}M}{} adalah

\mavergleichskettedisp
{\vergleichskette
{ \sqrt { \det  \begin{pmatrix}  1+v^2  &   -3v^3   \\  -3  v^3 &  1+ 9v^4  \end{pmatrix} }
}
{ =} { \sqrt{  (1+v^2)(1+9v^4) - 9v^6 }
}
{ =} { \sqrt{  1+v^2 +9v^4  }
}
{ } {
}
{ } {
}
}
{}{}{.}

}



\inputaufgabeklausurloesung
{6}

{

Buktikan aturan hasil kali untuk bentuk diferensial yang diferensiabel pada suatu himpunan terbuka di $\R^n$.

}
{

Misalkan
\mathl{x_1 , \ldots , x_n}{} adalah koordinat pada $\R^n$. Karena linearitas $d$ dan
multilinearitas hasil kali baji,
kita dapat menulis kedua bentuk diferensial sebagai
\mathkor {} {\omega=f dx_I} {dan} {\tau=g dx_J} {}
dengan himpunan indeks
\mathkor {} {I=\{i_1 , \ldots , i_k\}} {dan} {J=\{ j_1 , \ldots , j_\ell \}} {}.
Maka

\mavergleichskettealignhandlinks
{\vergleichskettealignhandlinks
{ d( \omega \wedge \tau)
}
{ =} { d \left(  fdx_I \wedge gdx_J  \right)
}
{ =} { d \left(  (f g) dx_I \wedge dx_J  \right)
}
{ =} { \sum_{s = 1}^n { \frac{   \partial  (fg)  }{  \partial   x_s  } }  dx_s  \wedge  dx_I  \wedge dx_J
}
{ =} { \sum_{s = 1}^n \left(  g { \frac{   \partial   f   }{  \partial   x_s  } } + f { \frac{   \partial   g   }{  \partial   x_s  } }  \right)  dx_s \wedge  dx_I  \wedge dx_J
}
}
{
\vergleichskettefortsetzungalign
{ =} { \sum_{s = 1}^n  g { \frac{   \partial   f   }{  \partial   x_s  } } dx_s \wedge  dx_I \wedge dx_J + \sum_{s = 1}^n  f { \frac{   \partial   g   }{  \partial   x_s  } }  dx_s \wedge  dx_I  \wedge dx_J
}
{ =} { \sum_{s = 1}^n  { \frac{   \partial   f   }{  \partial   x_s  } } dx_s \wedge  dx_I \wedge gdx_J + \sum_{s = 1}^n  { \frac{   \partial   g   }{  \partial   x_s  } }  dx_s \wedge f dx_I \wedge dx_J
}
{ =} { d(fdx_I ) \wedge g dx_J  + \sum_{s = 1}^n (-1)^k  f dx_I \wedge { \frac{   \partial   g   }{  \partial   x_s  } }  dx_s \wedge dx_J
}
{ =} { d(fdx_I ) \wedge g dx_J + (-1)^k  f dx_I \wedge  d( g dx_J)
}
}
{}{.}

}



\inputaufgabeklausurloesung
{0}

{

}
{  }



\inputaufgabeklausurloesung
{6}

{

Misalkan $M$ adalah suatu
\definitionsverweis {manifold berbatas}{}{}
dan misalkan

\mavergleichskette
{\vergleichskette
{ P
}
{ \in }{ \partial M
}
{  }{
}
{    }{
}
{   }{
}
}
{}{}{}
adalah titik batas. Tunjukkan bahwa sifat-sifat berikut ekuivalen bagi suatu vektor singgung

\mavergleichskette
{\vergleichskette
{ v
}
{ \in }{ T_PM
}
{  }{
}
{    }{
}
{   }{
}
}
{}{}{.}
\aufzaehlungzwei {Terdapat suatu lintasan yang terdiferensialkan secara kontinu
\maabb {\gamma} { [- \epsilon, 0]} { M
} {}
dengan

\mavergleichskette
{\vergleichskette
{ \gamma(0)
}
{ = }{ P
}
{  }{
}
{    }{
}
{   }{
}
}
{}{}{}
dan

\mavergleichskette
{\vergleichskette
{ \gamma'(0)
}
{ = }{ v
}
{  }{
}
{    }{
}
{   }{
}
}
{}{}{.}
} { Pada setiap peta bermodel setengah ruang negatif, $v$ dipetakan oleh pemetaan singgung ke setengah ruang kanan.
}

}
{

Misalkan $\gamma$ suatu setengah lintasan di $M$ seperti yang diberikan, dan misalkan

\mavergleichskette
{\vergleichskette
{ P
}
{ \in }{  U
}
{ \subseteq }{ M
}
{    }{
}
{   }{
}
}
{}{}{}
suatu daerah peta dengan peta
\maabbdisp {\alpha} { U } {V \subseteq  H_{\leq 0}
} {}
dengan

\mavergleichskettedisp
{\vergleichskette
{  H_{\leq 0}
}
{ =} { \left\{ (x_1,x_2 , \ldots, x_n)  \mid   x_1 \leq 0        \right\}
}
{ } {
}
{ } {
}
{ } {
}
}
{}{}{}
dan

\mavergleichskettedisp
{\vergleichskette
{ Q
}
{ =} { \alpha(P)
}
{ =} { \left(  0   , \,  a_2  , \,   \ldots , \,  a_n                \right)
}
{ } {
}
{ } {
}
}
{}{}{.}
Di bawah pemetaan tangen,

\mavergleichskette
{\vergleichskette
{ v
}
{ = }{ \gamma'(0)
}
{  }{
}
{    }{
}
{   }{
}
}
{}{}{}
dipetakan ke vektor turunan dari

\mavergleichskette
{\vergleichskette
{ \delta
}
{ = }{ \alpha \circ \gamma
}
{  }{
}
{    }{
}
{   }{
}
}
{}{}{}
di titik nol. Karena $\gamma$ berjalan dalam $M$, seluruh $\delta$ berjalan dalam $H_{\leq 0}$. Oleh karena itu, dalam hasil bagi selisih

\mathdisp {{ \frac{   \left(  \delta_1(t)   , \,   \delta_2(t)  , \,   \ldots , \,   \delta_n(t)                \right)  }{ t } }} {  }

baik $t$ maupun $\delta_1(t)$ bernilai negatif, sehingga

\mavergleichskettedisp
{\vergleichskette
{  \delta_1'(0)
}
{ \geq} { 0
}
{ } {
}
{ } {
}
{ } {
}
}
{}{}{,}
dan dengan demikian termasuk dalam setengah ruang positif.

Sebaliknya, misalkan
\mathl{T_P(\alpha)(v)}{} termasuk dalam setengah ruang positif. Maka garis
\maabbeledisp {\delta} {\R} {\R^n
} { t } { Q+tv
} {}
untuk

\mavergleichskette
{\vergleichskette
{ t
}
{ \leq }{  0
}
{  }{
}
{    }{
}
{   }{
}
}
{}{}{}
seluruhnya berjalan dalam setengah ruang negatif. Lintasan
\mathl{\alpha^{-1} \circ \delta}{} didefinisikan pada suatu interval negatif, bernilai dalam $M$, dan merupakan realisasi setengah lintasan dari $v$.

}



\inputaufgabeklausurloesung
{12}

{

Buktikan teorema Stokes untuk manifold berbatas.

}
{

\teilbeweis {}{}{}
{Misalkan

\mathbed {U_i} {}
 {i \in I} {}
 {} {} {} {,}
suatu
\definitionsverweis {penutup terbuka}{}{}
dari $M$ dengan
\definitionsverweis {peta-peta berorientasi}{}{},
dan misalkan

\mathbed {h_j} {}
 {j \in J} {}
 {} {} {} {,}
suatu
\definitionsverweis {partisi kesatuan yang terdiferensialkan secara kontinu dan subordinat terhadap penutup tersebut}{}{,}
yang ada menurut
Teorema 22.10 (Geometri diferensial (Osnabrück 2023)).
Untuk setiap

\mavergleichskette
{\vergleichskette
{ P
}
{ \in }{ M
}
{  }{
}
{    }{
}
{   }{
}
}
{}{}{}
terdapat lingkungan terbuka

\mavergleichskette
{\vergleichskette
{ P
}
{ \in }{ V_P
}
{ \subseteq }{ M
}
{    }{
}
{   }{
}
}
{}{}{}
sedemikian sehingga

\mavergleichskette
{\vergleichskette
{ h_j \vert_{V_P }
}
{ = }{ 0
}
{  }{
}
{    }{
}
{   }{
}
}
{}{}{}
kecuali untuk berhingga banyak indeks. Misalkan $Y$ dukungan dari $\omega$. Penutup

\mavergleichskette
{\vergleichskette
{ Y
}
{ \subseteq }{ \bigcup_{P \in M} V_P
}
{  }{
}
{    }{
}
{   }{
}
}
{}{}{}
memiliki subpenutup berhingga karena
\definitionsverweis {kekompakan}{}{}
yang diasumsikan, katakanlah

\mavergleichskettedisp
{\vergleichskette
{ Y
}
{ \subseteq} { V_{1} \cup \ldots \cup V_{r}
}
{ =} { V
}
{ } {
}
{ } {
}
}
{}{}{.}
Jadi, hanya berhingga banyak $h_j$ pada $V$ yang tidak sama dengan $0$. Kita tetapkan

\mavergleichskette
{\vergleichskette
{ \omega_j
}
{ = }{ h_j \omega
}
{  }{
}
{    }{
}
{   }{
}
}
{}{}{;}
bentuk-bentuk diferensial ini juga terdiferensialkan secara kontinu. Dukungan $h_j$ adalah suatu subhimpunan dalam $M$ yang
\definitionsverweis {tertutup}{}{,}
dan terletak dalam
\mathl{U_{i(j)}}{}; karena itu, dukungan $\omega_j$ terletak dalam
\mathl{Y \cap U_{i(j)}}{} dan bersifat kompak menurut
Soal 13.10 (Geometri diferensial (Osnabrück 2023)).
Berlaku

\mavergleichskettedisp
{\vergleichskette
{ \omega
}
{ =} { \sum_{j \in J} \omega_j
}
{ } {
}
{ } {
}
{ } {
}
}
{}{}{,}
dan hanya berhingga banyak bentuk diferensial
\mathl{\omega_j}{} ini yang tidak sama dengan $0$ karena

\mavergleichskette
{\vergleichskette
{ \omega_j \vert_{M \setminus Y}
}
{ = }{ 0
}
{  }{
}
{    }{
}
{   }{
}
}
{}{}{}
untuk semua $j$, sedangkan

\mavergleichskettedisp
{\vergleichskette
{ \omega_j \vert_V
}
{ =} { 0
}
{ } {
}
{ } {
}
{ } {
}
}
{}{}{}
untuk semua $j$ kecuali berhingga banyak. Berdasarkan
aditivitas integral bentuk diferensial
dan
aditivitas turunan eksterior,
pernyataan dapat dibuktikan secara terpisah untuk setiap $\omega_j$. Jadi, kita dapat mengasumsikan bahwa diberikan suatu bentuk diferensial berderajat $(n-1)$ yang terdiferensialkan secara kontinu, memiliki dukungan kompak, dan seluruh dukungannya terletak dalam suatu lingkungan peta $U$.}
{}
\teilbeweis {}{}{}
{Terdapat suatu
\definitionsverweis {difeomorfisme}{}{}
\maabb {\alpha} { U } { V
} {}
dengan

\mavergleichskette
{\vergleichskette
{ V
}
{ \subseteq }{ H
}
{  }{
}
{    }{
}
{   }{
}
}
{}{}{}
terbuka, yang sekaligus menginduksi difeomorfisme antara batas-batas

\mavergleichskette
{\vergleichskette
{ \partial U
}
{ = }{ \partial M \cap U
}
{  }{
}
{    }{
}
{   }{
}
}
{}{}{}
dan

\mavergleichskette
{\vergleichskette
{ \partial V
}
{ = }{ \partial H \cap V
}
{  }{
}
{    }{
}
{   }{
}
}
{}{}{}.
Di sini berlaku

\mavergleichskettedisp
{\vergleichskette
{
\int_{  \partial U  }   \omega
}
{ =} {
\int_{  \partial V  }  \alpha^{-1 *}\omega
}
{ } {
}
{ } {
}
{ } {
}
}
{}{}{}
dan

\mavergleichskettedisp
{\vergleichskette
{
\int_{ U  }  d \omega
}
{ =} {
\int_{ V  }  \alpha^{-1 *} (  d \omega)
}
{ =} {
\int_{  V  }  d \left(  \alpha^{-1 *}\omega  \right)
}
{ } {
}
{ } {
}
}
{}{}{}
menurut
Lema 20.2 (Geometri diferensial (Osnabrück 2023))  (5).}
{\leerzeichen{}Dengan demikian, kita dapat mulai dari suatu bentuk diferensial yang terdiferensialkan secara kontinu dan memiliki dukungan kompak, didefinisikan pada

\mavergleichskette
{\vergleichskette
{ V
}
{ \subseteq }{ H
}
{  }{
}
{    }{
}
{   }{
}
}
{}{}{},
yang dapat kita perluas menjadi bentuk nol pada seluruh $H$ di luar dukungannya.}
\teilbeweis {}{}{}
{Karena kekompakan dukungannya, terdapat suatu balok yang cukup besar

\mavergleichskette
{\vergleichskette
{ Q
}
{ \subseteq }{ H
}
{  }{
}
{    }{
}
{   }{
}
}
{}{}{,}
yang salah satu sisinya $S$ terletak pada $\partial H$ dan yang bertemu dukungan $\omega$ hanya di $S$. Pada semua sisi lain dari $Q$, $\omega$
\zusatzklammer {dan dengan demikian juga $d\omega$} {} {}
adalah bentuk nol. Oleh karena itu, di satu pihak

\mavergleichskettedisp
{\vergleichskette
{
\int_{ H  }  d\omega
}
{ =} {
\int_{ Q  }  d\omega
}
{ } {
}
{ } {
}
{ } {
}
}
{}{}{}
dan di pihak lain

\mavergleichskettedisp
{\vergleichskette
{
\int_{  \partial H  }   \omega
}
{ =} {
\int_{  S  }   \omega
}
{ =} {
\int_{  \partial Q  }   \omega
}
{ } {
}
{ } {
}
}
{}{}{.}
Jadi, pernyataan tersebut mengikuti dari versi balok teorema Stokes.}
{}

}



\inputaufgabeklausurloesung
{0}

{

}
{  }



\inputaufgabeklausurloesung
{9 (1+1+3+2+2)}

{

Tinjau parametrisasi trigonometrik
\maabbeledisp {\psi} {\R } { S^1
} { t } { \left(  \cos  t   , \,   \sin  t                   \right)
} {}
dari lingkaran satuan. Misalkan

\mavergleichskette
{\vergleichskette
{ c
}
{ \in }{ [0, 2 \pi[
}
{  }{
}
{    }{
}
{   }{
}
}
{}{}{}
dipilih tetap. Pada bundel vektor trivial
\maabbdisp {} {S^1 \times \R^2} { S^1
} {}
diberikan suatu
\definitionsverweis {koneksi linear}{}{}
$\nabla$ sedemikian sehingga, sepanjang $\psi$,
\definitionsverweis {koneksi tarik balik}{}{}
$\psi^* \nabla$ diberikan oleh
\definitionsverweis {simbol Christoffel}{}{}
\zusatzklammer {indeks bagi satu-satunya arah diferensiasi dihilangkan} {} {}

\mavergleichskettedisp
{\vergleichskette
{ \begin{pmatrix} \tilde{\Gamma}_1^1 (t)   &  \tilde{\Gamma}_1^2 (t)   \\ \tilde{\Gamma}_2^1(t)  &  \tilde{\Gamma}_2^2(t)  \end{pmatrix}
}
{ =} { \begin{pmatrix}  0  &   - { \frac{  c  }{  2 \pi   } }   \\  { \frac{  c  }{  2 \pi  } }  &  0  \end{pmatrix}
}
{ } {
}
{ } {
}
{ } {
}
}
{}{}{.}
Untuk suatu titik

\mavergleichskette
{\vergleichskette
{ Q
}
{ = }{ \begin{pmatrix}  a  \\b   \end{pmatrix}
}
{ \in }{ \R^2
}
{    }{
}
{   }{
}
}
{}{}{}
tinjau pemetaan
\maabbeledisp {\Psi_{Q}} {\R} { S^1 \times \R^2
} { t } { \left(  \cos  t   , \,   \sin  t  , \,   M(t)  \begin{pmatrix}  a  \\b   \end{pmatrix}                 \right)
} {}
dengan

\mavergleichskettedisp
{\vergleichskette
{ M(t)
}
{ =} { \begin{pmatrix}  \cos  { \frac{  ct  }{  2 \pi   } }   &   -  \sin  { \frac{  ct  }{  2 \pi   } }    \\  \sin  { \frac{  ct  }{  2 \pi  } }   &  \cos  { \frac{  ct  }{  2 \pi   } }  \end{pmatrix}
}
{ } {
}
{ } {
}
{ } {
}
}
{}{}{.}
\aufzaehlungfuenf{Tentukan $\Psi_Q(0)$.
}{Tentukan $\Psi_Q(2 \pi)$.
}{Tunjukkan bahwa $\Psi_Q$ adalah suatu
\definitionsverweis {pengangkatan horizontal}{}{}
sepanjang $\psi$.
}{Tunjukkan bahwa
\maabbeledisp {} {\R^2} { \R^2
} { Q } { \Psi_Q(2 \pi)
} {,}
merupakan
\definitionsverweis {isometri linear}{}{}
\zusatzklammer {titik pangkal $(1,0)$ tidak dicantumkan di sini} {} {.}
}{Tunjukkan bahwa
\maabbeledisp {} {\Z} { \operatorname{SL}_{  2  } \! \left(  \R  \right)
} { k } { \left(  Q \mapsto \Psi_Q(2k \pi )  \right)
} {,}
merupakan
\definitionsverweis {homomorfisme grup}{}{.}
}

}
{

\aufzaehlungfuenf{Kita peroleh

\mavergleichskettealign
{\vergleichskettealign
{ \Psi_Q(0)
}
{ =} { \left(  \cos  0   , \,   \sin  0  , \,   M(0) \begin{pmatrix}  a  \\b   \end{pmatrix}                 \right)
}
{ =} { \left(  1   , \,   0  , \,   \begin{pmatrix}   1   &   0   \\  0   &  1  \end{pmatrix}    \begin{pmatrix}  a  \\b   \end{pmatrix}                 \right)
}
{ =} { \left(  1   , \,   0  , \,     \begin{pmatrix}  a  \\b   \end{pmatrix}                 \right)
}
{ } {
}
}
{}
{}{.}
}{Kita peroleh

\mavergleichskettealign
{\vergleichskettealign
{ \Psi_Q(2 \pi )
}
{ =} { \left(  \cos  2 \pi   , \,   \sin  2 \pi  , \,   M(2 \pi) \begin{pmatrix}  a  \\b   \end{pmatrix}                 \right)
}
{ =} { \left(  1   , \,   0  , \,   \begin{pmatrix}  \cos  c   &   - \sin c    \\  \sin c   &   \cos  c     \end{pmatrix} \begin{pmatrix}  a  \\b   \end{pmatrix}                 \right)
}
{ =} { \left(  1   , \,   0  , \,     \begin{pmatrix} a \cos  c   - b \sin c   \\ a \sin  c   - b \cos c    \end{pmatrix}                 \right)
}
{ } {
}
}
{}
{}{.}
}{Kita peroleh

\mavergleichskettealign
{\vergleichskettealign
{  \Psi_Q(t)
}
{ =} { \left(  \cos  t    , \,   \sin  t  , \,   \begin{pmatrix}   \cos  { \frac{  ct  }{  2 \pi   } }   &   - \sin  { \frac{  ct  }{  2 \pi   } }    \\  \sin  { \frac{  ct  }{  2 \pi  } }   &   \cos  { \frac{  ct  }{  2 \pi   } }  \end{pmatrix}  \begin{pmatrix}  a  \\b   \end{pmatrix}                 \right)
}
{ =} { \left(  \cos  t   , \,   \sin  t  , \,   \begin{pmatrix}  a \cos  { \frac{  ct  }{  2 \pi  } }  - b \sin  { \frac{  ct  }{  2 \pi  } }   \\ a \sin  { \frac{  ct  }{  2 \pi  } }  +   b  \cos  { \frac{  ct  }{  2 \pi } }    \end{pmatrix}                 \right)
}
{ =} { \left(  \cos  t   , \,   \sin  t  , \,     \left(  a \cos  { \frac{  ct  }{  2 \pi  } }  - b \sin  { \frac{  ct  }{  2 \pi  } }  \right)   e_1 +  \left(  a \sin  { \frac{  ct  }{  2 \pi  } }  +   b  \cos  { \frac{  ct  }{  2 \pi } }  \right)   e_2                 \right)
}
{ } {
}
}
{}
{}{,}
khususnya, ini merupakan suatu pengangkatan. Turunan vertikal dari penampang ini dalam $\R^2$ di atas $\R$ dapat ditentukan menggunakan simbol-simbol Christoffel yang diberikan; kita peroleh

\mavergleichskettealigndrucklinks
{\vergleichskettealigndrucklinks
{ \left(  \psi^* \nabla  \right)_\partial  \left(    \left(  a \cos  { \frac{  ct  }{  2 \pi  } }  - b \sin  { \frac{  ct  }{  2 \pi  } }  \right)   e_1 +  \left(  a \sin  { \frac{  ct  }{  2 \pi  } }  +   b  \cos  { \frac{  ct  }{  2 \pi } }  \right)   e_2  \right)
}
{ =} {   \left(  a \cos  { \frac{  ct  }{  2 \pi  } }  - b \sin  { \frac{  ct  }{  2 \pi  } }  \right)   \left(  \psi^* \nabla  \right)_\partial \left(  e_1  \right)  + \partial    \left(  a \cos  { \frac{  ct  }{  2 \pi  } }  - b \sin  { \frac{  ct  }{  2 \pi  } }  \right)   e_1    +   \left(  a \sin  { \frac{  ct  }{  2 \pi  } }  +   b  \cos  { \frac{  ct  }{  2 \pi } }  \right)      \left(  \psi^* \nabla  \right)_\partial \left(  e_2  \right) + \partial  \left(  a \sin  { \frac{  ct  }{  2 \pi  } }  +   b  \cos  { \frac{  ct  }{  2 \pi } }  \right)   e_2
}
{ =} { -  \left(  a \cos  { \frac{  ct  }{  2 \pi  } }  - b \sin  { \frac{  ct  }{  2 \pi  } }  \right)  { \frac{   c  }{  2 \pi } }  e_2  +   { \frac{   c  }{  2 \pi } } \left(  -a \sin  { \frac{  ct  }{  2 \pi  } }  - b \cos  { \frac{  ct  }{  2 \pi  } }  \right)   e_1    +   \left(  a \sin  { \frac{  ct  }{  2 \pi  } }  +   b  \cos  { \frac{  ct  }{  2 \pi } }  \right) { \frac{   c  }{  2 \pi } }   e_1    +  { \frac{   c  }{  2 \pi } }   \left(  a \cos  { \frac{  ct  }{  2 \pi  } } -  b  \sin  { \frac{  ct  }{  2 \pi } }  \right)   e_2
}
{ =} { 0
}
{ } {
}
}
{}{}{,}
sehingga terdapat suatu
\definitionsverweis {pengangkatan horizontal}{}{}.
}{Pemetaan yang dimaksud adalah
\maabbeledisp {} {\R^2} { \R^2
} { \begin{pmatrix}  a  \\b   \end{pmatrix} } {    \begin{pmatrix}  \cos  c   &   - \sin  c   \\  \sin  c  &  \cos  c  \end{pmatrix} \begin{pmatrix}  a  \\b   \end{pmatrix}
} {,}
yang linear. Matriks yang merepresentasikannya ortogonal
\zusatzklammer {suatu matriks rotasi} {} {,}
sehingga diperoleh suatu isometri.
}{Pemetaan
\mathl{\Psi (2k \pi)}{} adalah pemetaan linear
\maabbeledisp {} {\R^2} { \R^2
} {\begin{pmatrix}  a  \\b   \end{pmatrix} } { \begin{pmatrix}  \cos kc   &   - \sin kc   \\  \sin kc  &  \cos kc  \end{pmatrix} \begin{pmatrix}  a  \\b   \end{pmatrix}
} {,}
yaitu rotasi sebesar sudut $kc$. Matriks rotasi untuk jumlah dua sudut adalah hasil kali matriks dari kedua rotasi tersebut
\zusatzklammer {lihat
Fakta *****} {} {,}
sehingga diperoleh suatu
\definitionsverweis {homomorfisme grup}{}{}.
}

}
